% binder-placement hero pothos-placeholder-001; path assets/draft-placeholder.txt; 3.30 x 3.30 in; 1:1 square; minimum 792 x 792 px, preferred 990 x 990 px; crop-to-fill; photograph the whole vine sharply with crop room.
% binder-details none
\begin{minipage}[t][3.30in][t]{3.36in}\vspace{0pt}
{\color{rust}\sffamily\bfseries\fontsize{9}{11}\selectfont DRAFT • PROVISIONAL IDENTITY}\par\vspace{8pt}
{\fontsize{28}{30}\selectfont Pothos}\par\vspace{7pt}
{\itshape\fontsize{15}{18}\selectfont Provisional \textit{Epipremnum aureum}}\par\vspace{10pt}
{\sffamily\fontsize{9.5}{12}\selectfont \textbf{Working example:} \textit{E. aureum} care. Reported common name only; this does not identify the user’s specimen or cultivar.}
\end{minipage}\hfill
\begin{minipage}[t][3.30in][t]{3.30in}\vspace{0pt}\AssetHero\end{minipage}\par
\AssetDetails\vspace{.05in}
\CareRow{LIGHT / EXPOSURE}{Moderate to bright light; avoid direct sun. Prolonged low light may reduce variegation or foliage brightness. Exposure/cultivar remain unknown. \textbf{[POT-PSU]}}{SOIL / SUBSTRATE}{Use loose, well-aerated houseplant medium in a draining pot. No retained source supports a universal ingredient ratio; this is a qualitative starting point. \textbf{[POT-NCSU; POT-WISC]}}
\CareRow{WATER}{When medium is dry, water thoroughly and let excess drain; never leave the pot standing in water. Check moisture, not a calendar. \textbf{[POT-WISC; POT-OUT]}}{TEMPERATURE / SEASON}{Keep warm and away from cold drafts. Follow actual growth and dry-down; calendar winter does not by itself require dormancy. \textbf{[POT-WISC]}}
\CareRow{FEEDING / MAINTENANCE}{Feed lightly during active growth; reduce or stop when growth slows. Wipe dusty leaves; prune vines just above a leaf to promote branching. \textbf{[POT-WISC]}}{PROPAGATION}{Cut a stem with a node and its bud—not a detached leaf alone. In water, submerge the rooting node, keep foliage above water, then pot after roots form. \textbf{[POT-NCSU; POT-PROP]}}
\CareRow{TROUBLESHOOTING / SAFETY}{Check moisture/exposure before diagnosing rot, droop, or yellowing. Keep from chewing pets; ASPCA lists pothos toxic to cats/dogs and notes oral irritation if eaten. \textbf{[POT-PSU; POT-ASPCA]}}{NATURAL HISTORY / TRIVIA}{Candidate \textit{E. aureum} is a climbing/trailing evergreen aroid native to the Society Islands; it climbs with aerial roots. \textbf{[POT-NCSU]}}
\vfill
{\color{sage}\rule{\linewidth}{.6pt}}\par\vspace{4pt}{\sffamily\fontsize{8}{10}\selectfont\textbf{SOURCES} POT-NCSU=NC State plant profile; POT-PROP=NC State “Propagation”; POT-WISC=UW–Madison “Pothos”; POT-PSU=Penn State “Pothos”; POT-ASPCA=“Devils Ivy”; POT-OUT=Iowa State outdoor guide. Full records: \texttt{binder/entries/pothos/sources.yaml}.\\\textbf{REVISION} v0.2 • 2026-09-17. Identity/photos/local conditions pending; draft example, not specimen-qualified.}
